% Styles et briques communs aux figures du \cref{ch:spdnet}.
% Les teintes reprennent celles de gfx/tikZ/riemann/style.tex pour que les
% figures des deux chapitres se lisent comme une meme famille.
% Le texte du memoire ne fait que 355pt (12,5 cm): toutes les figures de ce
% chapitre s'ecoulent donc de haut en bas, jamais de gauche a droite.
% Ce fichier peut etre \input plusieurs fois.
\definecolor{slayer}{RGB}{90,90,90}      % boite de couche, gris neutre
\definecolor{sspd}{RGB}{150,40,40}       % ce qui vit sur le cone SPD
\definecolor{sstiefel}{RGB}{40,70,140}   % ce qui vit sur Stiefel
\definecolor{seucl}{RGB}{50,110,120}     % ce qui vit dans un espace euclidien
\definecolor{sback}{RGB}{175,95,45}      % retropropagation
\tikzset{
  spdnet/.style={
    >={Stealth[length=5pt,width=4pt]},
    line join=round,
    font=\small,
    node distance=5mm,
  },
  % Boite de couche generique. La largeur est fixee pour que la colonne soit
  % alignee et que le texte respire.
  slayerbox/.style={
    draw=slayer!70, thick, rounded corners=2pt, align=center,
    inner sep=4pt, minimum height=9mm, minimum width=42mm, fill=white,
  },
  sstiefelbox/.style={slayerbox, draw=sstiefel!75, fill=sstiefel!7},
  splainbox/.style={slayerbox, fill=slayer!5},
  seuclbox/.style={slayerbox, draw=seucl!75, fill=seucl!7},
  smeanbox/.style={slayerbox, draw=sspd!75, fill=sspd!9},
  % Boite etroite, pour les colonnes doubles.
  snarrow/.style={slayerbox, minimum width=32mm},
  sarrow/.style={->, thick, slayer!85},
  sbackarrow/.style={->, thick, sback, dashed},
  sblock/.style={draw=slayer!45, thick, dashed, rounded corners=3pt, inner sep=5pt},
  slabel/.style={font=\footnotesize, inner sep=2pt},
  stiny/.style={font=\scriptsize, inner sep=1.5pt},
  % Accolade de regroupement a droite ou a gauche d'une colonne.
  sbrace/.style={decorate, decoration={brace, amplitude=4pt}, slayer!70},
}
